% ============================================================
%  CHAPITRE 1 — VUE D'ENSEMBLE DU PROJET
% ============================================================
\chapter{Vue d'ensemble du projet}

\section{Introduction}

Ce chapitre présente le cadre général du projet de maintenance prédictive avec intelligence
artificielle. Nous y exposons le rôle des agents intelligents dans la surveillance des
systèmes, les limites des approches existantes, les objectifs de notre solution
\textbf{PC Technician Assistant}, ainsi que la méthodologie de développement Agile adoptée
et le planning détaillé des sprints.

\section{Contexte et domaine d'application}

\subsection{Le parc informatique comme actif critique}

Dans les organisations modernes, le parc informatique constitue un actif stratégique dont
la disponibilité conditionne directement la productivité. Un technicien informatique gère
en moyenne entre 50 et 200 postes de travail, serveurs et équipements réseau. La surveillance
manuelle de l'état de santé de chaque machine est non seulement chronophage, mais aussi
insuffisante pour détecter les signes précoces de dégradation.

Les composants les plus sujets aux défaillances dans un parc informatique sont :
\begin{itemize}
  \item \textbf{Les disques durs (HDD/SSD)} : composants soumis à une usure progressive,
        responsables de la majorité des pertes de données non récupérables ;
  \item \textbf{Les processeurs (CPU)} : surchauffe due à une ventilation défaillante ou
        à une charge excessive prolongée, détectable via les capteurs de température ;
  \item \textbf{La mémoire vive (RAM)} : dégradation progressive pouvant entraîner des
        plantages aléatoires et des corruptions de données en mémoire ;
  \item \textbf{Les alimentations électriques} : vieillissement des condensateurs entraînant
        des instabilités de tension difficiles à détecter sans instrumentation.
\end{itemize}

\subsection{Enjeux économiques de la maintenance non planifiée}

Les coûts liés aux pannes informatiques non anticipées sont multidimensionnels. Au-delà du
coût direct du remplacement matériel (majoré de 30 à 50\% en cas d'achat en urgence), les
coûts indirects incluent la perte de productivité des utilisateurs, le temps de récupération
des données, et l'impact sur la continuité de service. La maintenance prédictive vise à
transformer ces coûts réactifs en investissements planifiés et optimisés.

\section{Rôle des agents intelligents dans la surveillance système}

\subsection{Architecture de l'agent de collecte}

Un agent intelligent est un programme autonome capable de percevoir son environnement,
de traiter les informations collectées et d'agir en conséquence. Dans notre système,
l'agent Python est déployé sur chaque machine surveillée et est composé de quatre modules
distincts :

\begin{itemize}
  \item \texttt{collector.py} : collecte les métriques CPU, RAM et disque via \texttt{psutil} ;
  \item \texttt{smart\_reader.py} : lit les attributs SMART du disque via \texttt{smartctl} ;
  \item \texttt{sender.py} : envoie les données au backend avec mécanisme de retry ;
  \item \texttt{scheduler.py} : orchestre la collecte horaire via la bibliothèque
        \texttt{schedule}.
\end{itemize}

\subsection{Métriques collectées par l'agent}

La classe \texttt{SystemCollector} collecte trois catégories de métriques :

\textbf{Métriques CPU :} l'usage est mesuré via \texttt{psutil.cpu\_percent(interval=1)},
qui effectue une mesure sur une seconde pour éviter les valeurs instantanées non
représentatives. La température est lue via \texttt{psutil.sensors\_temperatures()},
avec une recherche séquentielle des capteurs connus (\texttt{coretemp}, \texttt{cpu\_thermal},
\texttt{k10temp}, \texttt{zenpower}). En cas d'indisponibilité des capteurs (fréquent sur
les machines virtuelles), une valeur de fallback de 50°C est utilisée.

\textbf{Métriques mémoire :} \texttt{psutil.virtual\_memory()} retourne le pourcentage
d'utilisation et la mémoire disponible en mégaoctets, calculée par conversion depuis les
octets bruts.

\textbf{Métriques disque :} \texttt{psutil.disk\_usage()} est appelé sur le chemin
\texttt{C:\textbackslash\textbackslash} sous Windows ou \texttt{/} sous Linux, retournant
l'usage en pourcentage et l'espace libre en mégaoctets.

\subsection{Identification unique des machines}

L'identification de chaque machine repose sur une stratégie à trois niveaux de fallback :

\begin{enumerate}
  \item \textbf{Numéro de série BIOS} : via \texttt{wmic bios get serialnumber} (Windows)
        ou \texttt{/sys/class/dmi/id/product\_serial} (Linux). Les valeurs génériques
        (\texttt{TO BE FILLED BY O.E.M.}) sont rejetées ;
  \item \textbf{Adresse MAC} : utilisée si le numéro de série est invalide, formatée
        comme \texttt{MAC-xx:xx:xx:xx:xx:xx} ;
  \item \textbf{Identifiant UUID} : généré via \texttt{uuid.getnode()} comme dernier
        recours, préfixé \texttt{FALLBACK-}.
\end{enumerate}

\subsection{Mécanisme de résilience : retry avec backoff exponentiel}

La classe \texttt{DataSender} implémente un mécanisme de retry robuste pour gérer les
défaillances réseau transitoires. Le comportement est le suivant :

\begin{table}[H]
\centering
\caption{Comportement du mécanisme de retry de l'agent}
\begin{tabular}{|c|c|l|}
\hline
\rowcolor{lightblue}
\textbf{Tentative} & \textbf{Délai d'attente} & \textbf{Condition de déclenchement} \\
\hline
1 & 0s (immédiat) & Première tentative \\
\hline
2 & 1s & Erreur HTTP 5xx ou timeout réseau \\
\hline
3 & 2s & Erreur HTTP 5xx ou timeout réseau \\
\hline
— & Abandon & Erreur HTTP 4xx (erreur client) ou 3 tentatives épuisées \\
\hline
\end{tabular}
\end{table}

Les erreurs 4xx (Bad Request, Unauthorized) ne déclenchent pas de retry car elles indiquent
un problème de configuration (token invalide, payload malformé) qui ne se résoudra pas avec
une nouvelle tentative. Seules les erreurs 5xx (erreurs serveur) et les timeouts réseau
justifient un retry (\texttt{max\_retries = 3}, délais : 1s, 2s).

\section{Problématique et limites des systèmes existants}

\subsection{La maintenance réactive : un modèle dépassé}

La maintenance réactive consiste à intervenir uniquement après la survenue d'une panne.
Ce modèle présente plusieurs inconvénients majeurs :

\begin{itemize}
  \item Temps d'arrêt non planifiés impactant la productivité des utilisateurs ;
  \item Coûts élevés liés aux interventions d'urgence et aux achats de matériel en urgence ;
  \item Risque de perte de données en cas de défaillance soudaine d'un disque dur ;
  \item Impossibilité de planifier les ressources humaines et matérielles à l'avance.
\end{itemize}

\subsection{Limites des outils de monitoring existants}

Des outils comme Nagios, Zabbix ou Datadog offrent des capacités de surveillance en temps
réel basées sur des seuils fixes, mais ne disposent pas de capacités prédictives basées
sur l'apprentissage automatique :

\begin{table}[H]
\centering
\caption{Comparaison des approches de maintenance}
\begin{tabularx}{\textwidth}{|l|X|X|X|}
\hline
\rowcolor{lightblue}
\textbf{Critère} & \textbf{Réactive} & \textbf{Monitoring classique} & \textbf{Notre solution} \\
\hline
Détection & Après panne & Seuils fixes & Prédiction ML \\
\hline
Anticipation & Aucune & Limitée & 7/14/30 jours \\
\hline
ML intégré & Non & Non & Oui (RF + LSTM) \\
\hline
Coût & Élevé & Moyen (500€/mois) & Faible \\
\hline
Données SMART & Non & Partiel & Complet \\
\hline
Chatbot IA & Non & Non & Oui (LLaMA local) \\
\hline
Confidentialité & N/A & Cloud requis & 100\% local \\
\hline
\end{tabularx}
\end{table}

\section{Objectifs et solution proposée}

\subsection{Vision globale : PC Technician Assistant}

Notre solution repose sur quatre piliers fondamentaux :

\begin{enumerate}
  \item \textbf{Collecte automatisée} : un agent Python déployé sur chaque machine collecte
        les métriques système et SMART toutes les heures, de manière transparente ;
  \item \textbf{Prédiction par apprentissage automatique} : deux modèles complémentaires
        (Random Forest sur 63 features statistiques, LSTM sur séquences SMART Backblaze)
        prédisent les pannes à 7, 14 et 30 jours ;
  \item \textbf{Alertes intelligentes} : génération automatique d'alertes par email pour
        les risques HIGH ($\geq 50\%$) et CRITICAL ($\geq 70\%$) ;
  \item \textbf{Interface de gestion} : tableau de bord React avec visualisation des
        métriques, gestion des alertes, et chatbot conversationnel LLaMA local.
\end{enumerate}

\subsection{Niveaux de risque et seuils de décision}

Le système classifie chaque machine selon quatre niveaux de risque basés sur la probabilité
de panne à 30 jours calculée par le modèle Random Forest :

\begin{table}[H]
\centering
\caption{Classification des niveaux de risque et actions associées}
\begin{tabularx}{\textwidth}{|c|c|X|}
\hline
\rowcolor{lightblue}
\textbf{Probabilité (30j)} & \textbf{Niveau} & \textbf{Action déclenchée} \\
\hline
$< 30\%$ & LOW & Surveillance normale, aucune action \\
\hline
$30\% - 50\%$ & MEDIUM & Surveillance accrue, vérification planifiée \\
\hline
$50\% - 70\%$ & HIGH & Création d'alerte + envoi email au technicien \\
\hline
$\geq 70\%$ & CRITICAL & Alerte urgente + email prioritaire \\
\hline
\end{tabularx}
\end{table}

\section{Méthodologie de développement}

\subsection{Approche Agile en quatre sprints}

Le développement a suivi une approche Agile organisée en quatre sprints, chacun produisant
un incrément fonctionnel testable :

\begin{table}[H]
\centering
\caption{Planning des sprints de développement}
\begin{tabularx}{\textwidth}{|c|X|X|}
\hline
\rowcolor{lightblue}
\textbf{Sprint} & \textbf{Période} & \textbf{Livrables principaux} \\
\hline
Sprint 1 & 01 Fév -- 28 Fév 2026 & Agent Python, API Backend Node.js, Base de données PostgreSQL (9 tables), migrations SQL \\
\hline
Sprint 2 & 01 Mar -- 15 Mar 2026 & Dashboard React, Authentification JWT, RBAC Admin/Technicien, composants KPI \\
\hline
Sprint 3 & 16 Mar -- 31 Mar 2026 & Service ML Flask (RF + LSTM), Système d'alertes, Emails SMTP, APScheduler \\
\hline
Sprint 4 & 01 Avr -- 30 Avr 2026 & Chatbot Ollama/LLaMA, Évaluation ROUGE, Docker Compose, Finalisation \\
\hline
\end{tabularx}
\end{table}

\subsection{Gestion du backlog produit}

Le backlog produit a été organisé selon les priorités MoSCoW :

\begin{description}
  \item[\textbf{Must Have (MVP)}] : collecte des métriques, stockage PostgreSQL, prédictions
        Random Forest, alertes email, tableau de bord, authentification JWT ;
  \item[\textbf{Should Have}] : modèle LSTM Backblaze, chatbot IA, export CSV, gestion
        des utilisateurs, versioning des modèles ML ;
  \item[\textbf{Could Have}] : déploiement Docker Compose, évaluation ROUGE, détection
        d'anomalies Isolation Forest ;
  \item[\textbf{Won't Have (v1)}] : application mobile, module de ticketing, rapports PDF.
\end{description}

\section{Conclusion}

Ce chapitre a posé les bases conceptuelles du projet en définissant le problème adressé,
les objectifs visés, l'architecture de l'agent de collecte, et la méthodologie adoptée.
L'approche Agile en quatre sprints a permis de livrer un système complet et fonctionnel
dans les délais impartis. Le chapitre suivant présente l'état de l'art des technologies
utilisées, notamment les modèles d'apprentissage automatique pour la maintenance prédictive
et les approches d'évaluation des systèmes conversationnels.
